\begin{thebibliography}{99}
\bibitem{shigeno2020} 茂野仁美，乳幼児保育におけるリズムへの同期の発達過程に関する文献研究．大阪総合保育大学紀要，(14)，213-226，2019．
\bibitem{takasino} 高橋むつみ，篠田晴男，幼児期におけるリズム知覚と協応動作の発達：音楽活動を通した発達支
援の一助として，茨城大学教育学部紀要(教育科学)，(52)，249-262，2003．
\bibitem{Lazzaro} Lazzaro,J., \& Wawryzynek,J.(2001)."A case for network musical performance."Pro-
ceedings of the 11th international workshop on Network and operating systems support for
digital audio and video (NOSSDAV).
\bibitem{Rottondi} Rottondi, C., Chafe, C., Allocchio, C., \& Sart, A. (2016). ”An Overview on Networked
Music Performance Technologies.” IEEE Access, 4, 8823-8843.
\bibitem{NTPSNTP} D. Mills, J. Martin, J. Burbank and W. Kasch, “Network Time Protocol Version 4:
Protocol and Algorithms Specification,” RFC 5905, Jun. 2010.
\bibitem{PTP} IEEE, “IEEE Standard for a Precision Clock Synchronization Protocol for Networked
Measurement and Control Systems,” IEEE Std 1588-2019 (Revision of IEEE Std 1588-
2008), pp. 1–499, Nov. 2019.
\end{thebibliography}